\documentclass[10pt, aspectratio=169]{beamer}

% --- TECHNICAL CONTEXT & PACKAGE LOAD ---
\usepackage{amsmath}
\usepackage{booktabs}
\usepackage{listings}
\usepackage{xcolor}
\usepackage{fontspec}
\usepackage[english, bidi=basic, provide=*]{babel}

% --- FONT CONFIGURATION (Universal Preamble Block for Canvas LaTeX) ---
\babelfont{rm}{Noto Sans}
\babelfont{sf}{Noto Sans}

% --- BEAMER THEME & COLOR SCHEME ---
\usetheme{Boadilla}
\usecolortheme{whale}

% NSO Custom Corporate Colors
\definecolor{nsoblue}{RGB}{20, 50, 110}
\definecolor{nsogold}{RGB}{197, 160, 89}
\definecolor{darkgray}{RGB}{60, 60, 60}

\setbeamercolor{palette primary}{bg=nsoblue,fg=white}
\setbeamercolor{palette secondary}{bg=nsoblue!85,fg=white}
\setbeamercolor{palette tertiary}{bg=nsoblue!70,fg=white}
\setbeamercolor{titlelike}{parent=palette primary}
\setbeamercolor{structure}{fg=nsoblue}
\setbeamercolor{item}{fg=nsogold}
\setbeamercolor{block title}{bg=nsoblue,fg=white}
\setbeamercolor{block body}{bg=nsoblue!5,fg=black}
\setbeamercolor{block title alerted}{bg=red!80!black,fg=white}
\setbeamercolor{block body alerted}{bg=red!5,fg=black}

% --- LISTINGS CONFIGURATION FOR CODE ---
\lstset{
  basicstyle=\ttfamily\tiny,
  breaklines=true,
  backgroundcolor=\color{black!5},
  frame=single,
  rulecolor=\color{nsoblue!30},
  keywordstyle=\color{blue},
  commentstyle=\color{gray},
  stringstyle=\color{red!70!black},
  showstringspaces=false,
  escapeinside={(*}{*)}
}

% --- PRESENTATION METADATA ---
\title[Survey Weights \& Calibration]{Survey Weights, Post-Stratification, \\ and Weight Adjustment}
\subtitle{Session 2: Weights Theory and Calculation using Stata \& Excel}
\author[Dr. Anil Shrestha]{Dr. Anil Shrestha}
\institute[NSO]{Prepared for National Statistics Office (NSO) Training Participants \\ and Survey Practitioners}
\date[June 2026]{Issued: June 2026}

\begin{document}

% --- SLIDE 1: TITLE SLIDE ---
\begin{frame}[plain]
  \titlepage
\end{frame}

% --- SLIDE 2: ROADMAP ---
\begin{frame}{Learning Roadmap Overview}
  \begin{columns}[T]
    \begin{column}{0.5\textwidth}
      \begin{block}{Theoretical Concepts}
        \begin{itemize}
          \item \textbf{Step 1: Why Do We Need Weights?} \\ Core selection problems and sample imbalances.
          \item \textbf{Step 2: Design Weights} \\ Sampling probability foundations (SRS, Stratified, Two-Stage).
          \item \textbf{Step 3: Non-Response Adjustment} \\ Compensating for systemic field non-participation.
          \item \textbf{Step 4: Post-Stratification} \\ Correcting residual sample-population mismatches.
        \end{itemize}
      \end{block}
    \end{column}
    \begin{column}{0.5\textwidth}
      \begin{block}{Applications \& Diagnostics}
        \begin{itemize}
          \item \textbf{Step 5: Calibration (Raking)} \\ Marginal alignment using Iterative Proportional Fitting.
          \item \textbf{Step 6: Unified Workflows} \\ End-to-end processing pipelines in Stata and Excel.
          \item \textbf{Step 7: Quality Diagnostics} \\ Validation tables, weight trimming, and Kish's ESS.
        \end{itemize}
      \end{block}
    \end{column}
  \end{columns}
\end{frame}

% --- SECTION 1 ---
{
\setbeamercolor{background canvas}{bg=nsoblue}
\begin{frame}[plain, c]
  \centering
  \color{white}
  \Huge \textbf{Section 1: Why Do We Need Weights?} \\
  \vspace{0.4cm}
  \Large \textit{The Core Intuition and Sample Mismatch Problems}
\end{frame}
}

% --- SLIDE 3: CORE PROBLEM ---
\begin{frame}{The Core Problem: Unequal Selection Probabilities}
  \begin{itemize}
    \item \textbf{The Reality in National Surveys:} Not every population unit has an equal chance of being selected.
    \item \textbf{Geographical Representation:} To obtain reliable estimates in remote, sparsely populated regions (e.g., Humla), national agencies oversample them relative to highly populated urban centers (e.g., Kathmandu).
    \item \textbf{The Penalty of Unweighted Analysis:} If ignored, geographic and demographic groups are over-represented, causing severe bias in national indicators (e.g., poverty, average household income).
  \end{itemize}
  \begin{block}{Intuitive Definition: The Multiplication Factor}
    \text{A survey weight $w_i$ acts as a \textbf{multiplication factor}, indicating how many population units a single surveyed unit represents.}
    \begin{quote}
      \textit{``This household from Humla represents 175 households, while this Kathmandu household represents 600 households.''}
    \end{quote}
  \end{block}
\end{frame}

% --- SLIDE 4: SOURCES OF IMBALANCE ---
\begin{frame}{Sources of Sample Imbalance}
  There are three primary reasons why an achieved sample deviates from the target population:
  
  \begin{table}
    \centering
    \footnotesize
    \begin{tabular}{lll}
      \toprule
      \textbf{Source of Imbalance} & \textbf{What Happens in the Field} & \textbf{Methodological Remedy} \\
      \midrule
      Unequal Selection Probability & Some units are sampled more often by design & \textbf{Design weights} \\
      Non-response & Selected units refuse or cannot be reached & \textbf{Non-response adjustment} \\
      Sample-Population Mismatch & Achieved sample demographics differ from census & \textbf{Post-stratification / Raking} \\
      \bottomrule
    \end{tabular}
    \caption{Sources of Sample Imbalance and Remedies (Table 1)}
  \end{table}
\end{frame}

% --- SLIDE 5: THE WEIGHTED MEAN FORMULATION ---
\begin{frame}{The Weighted Mean Formulation}
  The weight $w_i$ attached to unit $i$ defines the number of elements in the target population represented by this specific respondent. If a household has a weight $w_i = 500$, it is treated as representing 500 households.
  
  \begin{block}{The Weighted Mean Equation}
    The weighted mean ($\overline{Y}_{w}$) is computed as:
    \begin{equation}
      \overline{Y}_{w} = \frac{\sum_{i=1}^{n} w_{i} \cdot Y_{i}}{\sum_{i=1}^{n} w_{i}}
    \end{equation}
    \begin{itemize}
      \item $Y_{i}$ is the surveyed value of interest for respondent $i$ (e.g., monthly income).
      \item $w_{i}$ is the final weight of respondent $i$.
      \item The denominator $\sum_{i=1}^{n} w_{i}$ sums all weights, which should approximate the total target population size ($N$).
    \end{itemize}
  \end{block}
\end{frame}

% --- SLIDE 6: WORKED EXAMPLE (DATA) ---
\begin{frame}{Worked Example: Household Sample Income Data}
  Consider a sample of $n=4$ households selected from two districts in Nepal (Table 2):
  
  \begin{table}
    \centering
    \small
    \begin{tabular}{llrr}
      \toprule
      \textbf{Household} & \textbf{District} & \textbf{Monthly Income ($Y_i$, NPR)} & \textbf{Weight ($w_i$)} \\
      \midrule
      Kathmandu A & Kathmandu & 80,000 & 120 \\
      Kathmandu B & Kathmandu & 60,000 & 120 \\
      Humla C & Humla & 20,000 & 850 \\
      Humla D & Humla & 25,000 & 850 \\
      \bottomrule
    \end{tabular}
    \caption{Stylized Sample Dataset (Table 2)}
  \end{table}

  \begin{block}{Unweighted Mean (Methodologically Incorrect)}
    Ignoring selection weights treats every household identically:
    \begin{equation*}
      \overline{Y}_{unweighted} = \frac{80,000 + 60,000 + 20,000 + 25,000}{4} = \mathbf{46,250\text{ NPR}}
    \end{equation*}
  \end{block}
\end{frame}

% --- SLIDE 7: WORKED EXAMPLE (WEIGHTED CALCULATION) ---
\begin{frame}{Worked Example: Weighted Mean Calculation}
  Applying weights using the proper weighted mean formulation:
  
  \begin{equation*}
    \overline{Y}_{w} = \frac{(120 \times 80,000) + (120 \times 60,000) + (850 \times 20,000) + (850 \times 25,000)}{120 + 120 + 850 + 850}
  \end{equation*}
  
  \begin{equation*}
    \overline{Y}_{w} = \frac{9,600,000 + 7,200,000 + 17,000,000 + 21,250,000}{1,940}
  \end{equation*}
  
  \begin{equation*}
    \overline{Y}_{w} = \frac{55,050,000}{1,940} \approx \mathbf{28,376.29\text{ NPR}}
  \end{equation*}

  \begin{block}{The Difference}
    The unweighted mean is severely biased upward (\textbf{46,250 NPR}) because it over-represents high-income urban Kathmandu households. The weighted mean (\textbf{28,376.29 NPR}) correctly represents the rural-dominant population profile.
  \end{block}
\end{frame}

% --- SLIDE 8: POLICY IMPLICATIONS ---
\begin{frame}{Policy Implications of Unweighted Estimations}
  \begin{itemize}
    \item \textbf{Substantial Bias:} The unweighted mean is biased upward by \textbf{63\%} (46,250 NPR vs. 28,376 NPR) because it over-represents urban Kathmandu.
    \item \textbf{Policy Failure:} Publishing the unweighted figure would cause CBS to overestimate national living standards. Social protection thresholds, budget allocations, and poverty interventions would target incorrect thresholds.
  \end{itemize}

  \begin{alertblock}{Key Thing to Hold Onto}
    A survey weight tells you how many population units a sampled unit represents. Weighted estimates correct for design-driven unequal selection, non-response, and residual mismatches. Never analyze representative survey datasets without applying weights.
  \end{alertblock}
\end{frame}

% --- SECTION 2 ---
{
\setbeamercolor{background canvas}{bg=nsoblue}
\begin{frame}[plain, c]
  \centering
  \color{white}
  \Huge \textbf{Section 2: Design Weights} \\
  \vspace{0.4cm}
  \Large \textit{Moving from Sampling Probability to Basic Weights}
\end{frame}
}

% --- SLIDE 9: FUNDAMENTAL RULES & SRS ---
\begin{frame}{Fundamental Rules and Simple Random Sampling (SRS)}
  \begin{block}{The Fundamental Reciprocal Rule}
    The design weight ($w_{i}^{design}$) is the exact mathematical inverse (reciprocal) of the inclusion probability ($\pi_{i}$):
    \begin{equation}
      w_{i}^{design} = \frac{1}{\pi_{i}}
    \end{equation}
    If a household has a 1-in-500 chance of selection, its design weight is exactly 500.
  \end{block}

  \begin{block}{Design 1: Simple Random Sampling (SRS)}
    In an SRS, every unit in the population has an equal, non-zero probability of selection:
    \begin{equation}
      \pi_{i} = \frac{n}{N} \Rightarrow w_{i}^{design} = \frac{N}{n}
    \end{equation}
    If Nepal has $N = 6,000,000$ households and you sample $n = 12,000$ households via SRS, all design weights are uniform: $w_i = 500$.
  \end{block}
\end{frame}

% --- SLIDE 10: STRATIFIED RANDOM SAMPLING ---
\begin{frame}{Design 2: Stratified Random Sampling}
  The population is divided into mutually exclusive subpopulations called strata ($h$), and an independent SRS is drawn from each stratum:
  \begin{equation}
    \pi_{ih} = \frac{n_h}{N_h} \Rightarrow w_{ih} = \frac{N_h}{n_h}
  \end{equation}

  \begin{table}
    \centering
    \footnotesize
    \begin{tabular}{lrrr}
      \toprule
      \textbf{Stratum ($h$)} & \textbf{Population ($N_h$)} & \textbf{Sample Size ($n_h$)} & \textbf{Design Weight ($w_{ih} = N_h/n_h$)} \\
      \midrule
      Urban & 1,800,000 & 3,000 & \textbf{600.00} \\
      Rural Hills & 2,500,000 & 5,000 & \textbf{500.00} \\
      Rural Mountains & 700,000 & 4,000 & \textbf{175.00} \\
      \midrule
      \textbf{Total} & \textbf{5,000,000} & \textbf{12,000} & \\
      \bottomrule
    \end{tabular}
    \caption{Stratified Sampling Design Weights (Table 3)}
  \end{table}

  \textbf{Interpretation:} Rural Mountains are oversampled ($0.57\%$) compared to Urban areas ($0.17\%$) to obtain stable regional estimates. The design weights correct for this, dropping from 600 (Urban) to 175 (Rural Mountains).
\end{frame}

% --- SLIDE 11: TWO-STAGE CLUSTER SAMPLING ---
\begin{frame}{Design 3: Two-Stage Cluster Sampling}
  The standard design utilized in large-scale household surveys (NLSS, DHS, MICS, etc.):
  \begin{itemize}
    \item \textbf{Stage 1:} Select $m$ Primary Sampling Units (PSUs, e.g., wards) with \textbf{Probability Proportional to Size (PPS)}.
    \item \textbf{Stage 2:} Select $n_{PSU}$ households within each selected PSU using \textbf{SRS}.
  \end{itemize}

  \begin{columns}[T]
    \begin{column}{0.48\textwidth}
      \begin{block}{Selection Probabilities}
        $$\pi_{PSU} = \frac{m \times M_{PSU}}{M_{total}}$$
        $$\pi_{HH|PSU} = \frac{n_{PSU}}{M_{PSU}}$$
        \textit{Where $M_{PSU}$ is the size measure (census household count) of the PSU, and $M_{total}$ is the stratum count.}
      \end{block}
    \end{column}
    \begin{column}{0.5\textwidth}
      \begin{block}{Combined Design Weight}
        $$\pi_i = \pi_{PSU} \times \pi_{HH|PSU} = \frac{m \times n_{PSU}}{M_{total}}$$
        \begin{equation}
          w_i = \frac{1}{\pi_i} = \frac{M_{total}}{m \times n_{PSU}}
        \end{equation}
        \textbf{Note:} The PSU-specific size measure $M_{PSU}$ completely cancels out of the joint weight equation!
      \end{block}
    \end{column}
  \end{columns}
\end{frame}

% --- SLIDE 12: TWO-STAGE WORKED EXAMPLE ---
\begin{frame}{Worked Example: Two-Stage Cluster Parameters}
  Let's calculate the design weight for a respondent household under Table 4 specifications:
  \begin{itemize}
    \item Total households in Stratum ($M_{total}$): \textbf{500,000}
    \item Number of PSUs selected ($m$): \textbf{50}
    \item Households in Specific PSU ($M_{PSU}$): \textbf{200}
    \item Households sampled within selected PSU ($n_{PSU}$): \textbf{20}
  \end{itemize}

  \begin{block}{Inclusion \& Weight Derivation}
    \begin{itemize}
      \item \textbf{Stage 1 Selection Probability:} $\pi_{PSU} = \frac{50 \times 200}{500,000} = 0.02$
      \item \textbf{Stage 2 Selection Probability:} $\pi_{HH|PSU} = \frac{20}{200} = 0.10$
      \item \textbf{Overall Inclusion Probability:} $\pi_i = 0.02 \times 0.10 = 0.002$
      \item \textbf{Final Design Weight:} $w_{i}^{design} = \frac{1}{0.002} = \mathbf{500}$
    \end{itemize}
  \end{block}
\end{frame}

% --- SLIDE 13: DESIGN WEIGHT IMPLEMENTATION (STATA & EXCEL) ---
\begin{frame}[fragile]{Constructing Design Weights: Stata and Excel}
  \begin{columns}[T]
    \begin{column}{0.55\textwidth}
      \begin{block}{Stata Script (Listing 1)}
\begin{lstlisting}[language=SQL]
* Compute overall probability
gen pi_i = prob_psu * prob_hh_given_psu

* Compute design weight
gen wt_design = 1 / pi_i

* Sanity check (sum of weights should approximate N)
summarize wt_design
display "Sum of weights: " r(sum)

* Declare complex survey structure to Stata
svyset psu [pweight = wt_design], strata(stratum) vce(linearized)

* Run survey-weighted mean estimation
svy: mean income
\end{lstlisting}
      \end{block}
    \end{column}
    \begin{column}{0.42\textwidth}
      \begin{block}{Excel Formula Setup}
        \textbf{Stratified Sampling:}
        \begin{itemize}
          \item Column E: \texttt{=C2/D2} \\ \textit{(i.e., Population\_Size / Sample\_Size)}
        \end{itemize}
        \textbf{Two-Stage Cluster:}
        \begin{itemize}
          \item Column D: \texttt{=B2*C2} \\ \textit{($\pi_i = \pi_{PSU} \times \pi_{HH|PSU}$)}
          \item Column E: \texttt{=1/D2} \\ \textit{($w_i^{design} = 1/\pi_i$)}
        \end{itemize}
        \textbf{Weighted Mean in Excel:}
        \begin{lstlisting}[backgroundcolor=\color{white}, frame=none, basicstyle=\ttfamily\scriptsize]
=SUMPRODUCT(weight_range, income_range) 
 / SUM(weight_range)
        \end{lstlisting}
      \end{block}
    \end{column}
  \end{columns}
\end{frame}

% --- SECTION 3 ---
{
\setbeamercolor{background canvas}{bg=nsoblue}
\begin{frame}[plain, c]
  \centering
  \color{white}
  \Huge \textbf{Section 3: Non-Response Adjustment} \\
  \vspace{0.4cm}
  \Large \textit{Handling Systemic and Non-Random Field Non-Participation}
\end{frame}
}

% --- SLIDE 14: THE THREAT ---
\begin{frame}{The Threat of Non-Random Non-Response}
  \begin{itemize}
    \item \textbf{The Ideal vs. Reality:} Design weights assume perfect participation. In the field, selected units fail to participate due to refusals, structural absences, or communication barriers.
    \item \textbf{Systemic Non-Response Bias:} If non-response is random (Missing Completely at Random, or MCAR), it merely drops sample size. If non-random (wealthier urban families refuse; highly mobile households are unreachable), it introduces severe bias.
    \item \textbf{Methodological Remedy:} Inflate the design weights of cooperating respondents within homogeneous \textbf{adjustment cells} to represent non-respondents.
  \end{itemize}

  \begin{columns}[T]
    \begin{column}{0.48\textwidth}
      \begin{block}{NR Adjustment Factor ($f_c^{NR}$)}
        \begin{equation}
          f_c^{NR} = \frac{\text{Sampled Units in Cell } c}{\text{Responding Units in Cell } c} = \frac{n_c}{r_c}
        \end{equation}
      \end{block}
    \end{column}
    \begin{column}{0.48\textwidth}
      \begin{block}{Adjusted Design Weight}
        \begin{equation}
          w_i^{adjusted} = w_i^{design} \times f_c^{NR}
        \end{equation}
      \end{block}
    \end{column}
  \end{columns}
\end{frame}

% --- SLIDE 15: ADJUSTMENT CELLS & EXAMPLES ---
\begin{frame}{Defining Non-Response Adjustment Cells}
  Cells must be defined using auxiliary variables known for both respondents and non-respondents:
  \begin{itemize}
    \item \textbf{Variables:} Stratum (Urban/Rural), Province, PSU Type, Household Size (larger families are easier to find at home), and Interviewer ID.
  \end{itemize}

  \begin{table}
    \centering
    \footnotesize
    \begin{tabular}{lrrrr}
      \toprule
      \textbf{Adjustment Cell ($c$)} & \textbf{Sampled ($n_c$)} & \textbf{Responded ($r_c$)} & \textbf{Response Rate} & \textbf{NR Factor ($n_c/r_c$)} \\
      \midrule
      Urban (Cell 1) & 3,000 & 2,700 & 90.0\% & \textbf{1.111} \\
      Rural Hills (Cell 2) & 5,000 & 4,600 & 92.0\% & \textbf{1.087} \\
      Rural Mountains (Cell 3) & 4,000 & 3,200 & 80.0\% & \textbf{1.250} \\
      \bottomrule
    \end{tabular}
    \caption{Field Outcomes and Non-Response Factors (Table 6)}
  \end{table}

  \begin{alertblock}{Cell Sizing \& MAR Assumption}
    Each non-response adjustment cell \textbf{must contain at least 20-30 respondents} to avoid highly unstable inflation factors. Adjustment relies on the \textbf{Missing at Random (MAR)} assumption: respondents and non-respondents share similar characteristics within the cell. Keep factors \textbf{below 2.0} to protect MAR.
  \end{alertblock}
\end{frame}

% --- SLIDE 16: NON-RESPONSE IMPLEMENTATION (STATA & EXCEL) ---
\begin{frame}[fragile]{Executing Non-Response Adjustments in Stata \& Excel}
  \begin{columns}[T]
    \begin{column}{0.55\textwidth}
      \begin{block}{Stata Implementation (Listing 4)}
\begin{lstlisting}[language=SQL]
* Method 1: Cell-based non-response adjustment
bysort stratum: egen n_sampled = count(hh_id)
bysort stratum: egen n_responded = total(responded)
gen nr_factor = n_sampled / n_responded
gen wt_nr_adjusted = wt_design * nr_factor

* EXCLUDE non-respondents for final analyses
keep if responded == 1

* Method 2: Response Propensity Modeling (Advanced)
logit responded i.stratum hh_size i.province [pweight=wt_design]
predict p_respond, pr
gen nr_factor_ps = 1 / p_respond
gen wt_propensity = wt_design * nr_factor_ps
\end{lstlisting}
      \end{block}
    \end{column}
    \begin{column}{0.42\textwidth}
      \begin{block}{Excel Lookup Setup}
        \textbf{Sheet 1: \texttt{Cell\_Factors}}\\
        Contains \texttt{Stratum}, \texttt{n\_sampled}, \texttt{n\_responded}, and factor formulas:
        \texttt{=B2/C2}
        
        \vspace{0.3cm}
        \textbf{Sheet 2: \texttt{HH\_Data}}\\
        Retrieve cell-based factor using \texttt{VLOOKUP}:
        \begin{lstlisting}[backgroundcolor=\color{white}, frame=none, basicstyle=\ttfamily\scriptsize]
Column D (NR_Factor):
=VLOOKUP(B2, Cell_Factors!$A:$D, 
         4, FALSE)

Column E (wt_adjusted):
=C2 * D2  '(wt_design * NR_Factor)
        \end{lstlisting}
      \end{block}
    \end{column}
  \end{columns}
\end{frame}

% --- SECTION 4 ---
{
\setbeamercolor{background canvas}{bg=nsoblue}
\begin{frame}[plain, c]
  \centering
  \color{white}
  \Huge \textbf{Section 4: Post-Stratification} \\
  \vspace{0.4cm}
  \Large \textit{Anchoring Survey Results to Known Census Distributions}
\end{frame}
}

% --- SLIDE 17: RESIDUAL MISMATCH ---
\begin{frame}{The Residual Mismatch Problem \& Post-Stratification}
  \begin{itemize}
    \item \textbf{Why does mismatch persist?} Even after design and non-response adjustments, survey estimates can still deviate from census demographics due to:
    \begin{itemize}
      \item \textbf{Sampling Variability:} Random fluctuations in who was selected.
      \item \textbf{Sampling Frame Imperfections:} Outdated registers and housing lists.
      \item \textbf{Residual Non-Response Bias:} Non-participation not covered by basic cells.
    \end{itemize}
    \item \textbf{Solution:} Scale weights within post-strata (e.g., Sex, Age, Region) to force the survey total to match external control benchmarks exactly (Table 7: Census/Administrative Registers).
  \end{itemize}

  \begin{columns}[T]
    \begin{column}{0.5\textwidth}
      \begin{block}{Post-Stratification Formula}
        \begin{equation}
          w_i^{PS} = w_i^{adjusted} \times \frac{N_g^{census}}{\hat{N}_g^{survey}}
        \end{equation}
      \end{block}
    \end{column}
    \begin{column}{0.46\textwidth}
      \begin{block}{Estimated Group Population}
        \begin{equation}
          \hat{N}_g^{survey} = \sum_{i \in g} w_i^{adjusted}
        \end{equation}
      \end{block}
    \end{column}
  \end{columns}
\end{frame}

% --- SLIDE 18: POST-STRATIFICATION WORKED EXAMPLE ---
\begin{frame}{Post-Stratification: Worked Example \& Factors}
  Consider a sample aligned to four post-strata based on Sex and Area (Table 8):
  
  \begin{table}
    \centering
    \footnotesize
    \begin{tabular}{llrrr}
      \toprule
      \textbf{Post-Stratum Group ($g$)} & \textbf{Demographics} & \textbf{Census ($N_g^{census}$)} & \textbf{Survey Total ($\hat{N}_g^{survey}$)} & \textbf{PS Factor} \\
      \midrule
      Group 1 & Male, Urban & 1,850,000 & 1,780,000 & \textbf{1.0393} \\
      Group 2 & Female, Urban & 1,920,000 & 1,970,000 & \textbf{0.9746} \\
      Group 3 & Male, Rural & 4,100,000 & 4,180,000 & \textbf{0.9808} \\
      Group 4 & Female, Rural & 4,300,000 & 4,240,000 & \textbf{1.0142} \\
      \bottomrule
    \end{tabular}
  \end{table}

  \begin{columns}[T]
    \begin{column}{0.48\textwidth}
      \begin{block}{Case A: Underrepresented Group (G1)}
        Household with $w_i^{adjusted} = 666.7$ is scaled up:
        $$w_i^{PS} = 666.7 \times 1.0393 = \mathbf{692.9}$$
      \end{block}
    \end{column}
    \begin{column}{0.48\textwidth}
      \begin{block}{Case B: Overrepresented Group (G2)}
        Household with $w_i^{adjusted} = 666.7$ is scaled down:
        $$w_i^{PS} = 666.7 \times 0.9746 = \mathbf{649.7}$$
      \end{block}
    \end{column}
  \end{columns}
\end{frame}

% --- SLIDE 19: POST-STRATIFICATION IMPLEMENTATION ---
\begin{frame}[fragile]{Implementing Post-Stratification: Stata and Excel}
  \begin{columns}[T]
    \begin{column}{0.55\textwidth}
      \begin{block}{Stata Implementation (Listing 7)}
\begin{lstlisting}[language=SQL]
* Method A: Manual Factor Calculation
bysort ps_group: egen survey_total = total(wt_nr)

* Assign Census Control Benchmarks
gen census_pop = .
replace census_pop = 1850000 if ps_group == 1
replace census_pop = 1920000 if ps_group == 2
replace census_pop = 4100000 if ps_group == 3
replace census_pop = 4300000 if ps_group == 4

* Apply factors
gen ps_factor = census_pop / survey_total
gen wt_poststrat = wt_nr * ps_factor

* Method B: Using Stata's Built-In svyset Engine
svyset psu [pweight = wt_nr], strata(stratum) ///
  poststrata(ps_group) postweight(census_pop) vce(linearized)
\end{lstlisting}
      \end{block}
    \end{column}
    \begin{column}{0.42\textwidth}
      \begin{block}{Excel Calculations}
        \textbf{Sheet: \texttt{PS\_Benchmarks}}\\
        Formula in Col E (Survey Weighted Total):
        \begin{lstlisting}[backgroundcolor=\color{white}, frame=none, basicstyle=\ttfamily\tiny]
=SUMIF(HH_Data!K:K, A2, HH_Data!J:J)
        \end{lstlisting}
        Col F (PS Factor): \texttt{=B2/E2}\\
        
        \vspace{0.25cm}
        \textbf{Sheet: \texttt{HH\_Data}}\\
        Retrieve factor from Col L:
        \begin{lstlisting}[backgroundcolor=\color{white}, frame=none, basicstyle=\ttfamily\tiny]
=VLOOKUP(K2, PS_Benchmarks!$A:$F, 
         6, FALSE)
        \end{lstlisting}
        Col M (wt\_final): \texttt{=J2*L2}
      \end{block}
    \end{column}
  \end{columns}
\end{frame}

% --- SECTION 5 ---
{
\setbeamercolor{background canvas}{bg=nsoblue}
\begin{frame}[plain, c]
  \centering
  \color{white}
  \Huge \textbf{Section 5: Calibration (Raking)} \\
  \vspace{0.4cm}
  \Large \textit{Overcoming the Dimensionality Trap with Iterative Proportional Fitting}
\end{frame}
}

% --- SLIDE 20: DIMENSIONALITY TRAP ---
\begin{frame}{The Multi-Dimensional Dimensionality Trap}
  \begin{itemize}
    \item \textbf{The Post-Stratification Constraint:} If we want to align the sample to multiple variables simultaneously (e.g., Province [7 categories], Sex [2], Age [6], and Area [2]), cross-classifying them yields:
    $$7 \times 2 \times 6 \times 2 = 168 \text{ joint cells}$$
    Many cells will have too few or zero respondents, producing unstable or mathematically impossible factors.
    \item \textbf{The Raking Solution:} Iterative Proportional Fitting (IPF) sequentially aligns weights to match multiple \textbf{marginal benchmarks} rather than the full cross-classified joint distribution.
  \end{itemize}

  \begin{block}{The Iterative Raking Loop}
    \centering
    Adjust to Marginal Benchmark 1 $\rightarrow$ Adjust to Marginal Benchmark 2 $\rightarrow$ Re-adjust to Benchmark 1 $\dots$ \\
    \textit{This repeats until weights converge and all marginal constraints are simultaneously satisfied.}
  \end{block}
\end{frame}

% --- SLIDE 21: RAKING WORKED EXAMPLE ---
\begin{frame}{Raking: Worked Numerical Example}
  Goal: Calibrate starting weights to two margins: Sex and Area. Total Population $N = 12,170,000$.
  
  \begin{columns}[T]
    \begin{column}{0.48\textwidth}
      \begin{block}{Marginal Census Benchmarks (Table 9)}
        \begin{tabular}{lr}
          \textbf{Sex Group} & \textbf{Census Count} \\
          \midrule
          Male & 5,950,000 \\
          Female & 6,220,000 \\
          \midrule
          \textbf{Area Group} & \textbf{Census Count} \\
          \midrule
          Urban & 3,770,000 \\
          Rural & 8,400,000 \\
        \end{tabular}
      \end{block}
    \end{column}
    \begin{column}{0.48\textwidth}
      \begin{block}{Pre-Raking Survey Totals (Table 10)}
        \begin{tabular}{lrrr}
          \toprule
          & \textbf{Urban} & \textbf{Rural} & \textbf{Total} \\
          \midrule
          \textbf{M} & 1,780,000 & 4,180,000 & 5,960,000 \\
          \textbf{F} & 1,970,000 & 4,240,000 & 6,210,000 \\
          \midrule
          \textbf{Total} & 3,750,000 & 8,420,000 & \textbf{12,170,000} \\
          \bottomrule
        \end{tabular}
      \end{block}
    \end{column}
  \end{columns}

  \begin{block}{Sequential Adjustments}
    \begin{itemize}
      \item \textbf{Iteration 1, Pass 1 (Sex):} Factor$_{Male}$ = 5.95M/5.96M $\approx$ 0.99832; Factor$_{Female}$ = 6.22M/6.21M $\approx$ 1.00161.
      \item \textbf{Iteration 1, Pass 2 (Area):} Recalculate margins with intermediate weights; adjust by Area Factors (Urban $\approx$ 1.00453; Rural $\approx$ 0.99798).
    \end{itemize}
  \end{block}
\end{frame}

% --- SLIDE 22: RAKING STATA CODE ---
\begin{frame}[fragile]{Implementing Calibration Raking in Stata}
  \begin{columns}[T]
    \begin{column}{0.48\textwidth}
      \begin{block}{Method A: User Package \texttt{ipfraking} (Listing 10)}
\begin{lstlisting}[language=SQL]
* ssc install ipfraking

ipfraking [pw = wt_nr], ///
  ctotal(sex: 5950000 6220000) ///
  ctotal(area: 3770000 8400000) ///
  generate(wt_raked) ///
  maxiter(50) tolerance(0.0001)

* Declare design with final raked weights
svyset psu [pweight = wt_raked], ///
  strata(stratum) vce(linearized)
\end{lstlisting}
      \end{block}
    \end{column}
    \begin{column}{0.48\textwidth}
      \begin{block}{The Convergence Criterion}
        The iteration sequence stops once the maximum relative difference between survey estimates and census benchmarks falls below threshold $\epsilon$ (typically $\epsilon = 0.0001$ or $0.01\%$):
        \begin{equation}
          \max_{g} \left| \frac{\hat{N}_g^{survey} - N_g^{census}}{N_g^{census}} \right| < \epsilon
        \end{equation}
      \end{block}
    \end{column}
  \end{columns}
\end{frame}

% --- SLIDE 23: MANUAL RAKING LOOP IN STATA (PART 1) ---
\begin{frame}[fragile]{Manual Programmatic Raking Loop in Stata (Part 1)}
  Below is the first part of the manual loop initialization and sequential pass calculation:
  \begin{lstlisting}[language=SQL]
gen wt_raked = wt_nr
local male_pop = 5950000; local female_pop = 6220000
local urban_pop = 3770000; local rural_pop = 8400000

local converged = 0; local iter = 0
while `converged' == 0 & `iter' < 50 {
    local iter = `iter' + 1
    
    * Pass 1: Sex margin adjustment
    bysort sex: egen s_sex = total(wt_raked)
    gen f_sex = cond(sex==1, `male_pop'/s_sex, `female_pop'/s_sex)
    replace wt_raked = wt_raked * f_sex
    drop s_sex f_sex
    
    * Pass 2: Area margin adjustment
    bysort area: egen s_area = total(wt_raked)
    gen f_area = cond(area==1, `urban_pop'/s_area, `rural_pop'/s_area)
    replace wt_raked = wt_raked * f_area
    drop s_area f_area
\end{lstlisting}
\end{frame}

% --- SLIDE 24: MANUAL RAKING LOOP IN STATA (PART 2) ---
\begin{frame}[fragile]{Manual Programmatic Raking Loop in Stata (Part 2)}
  This continues the manual loop, checking relative deviations to determine convergence:
  \begin{lstlisting}[language=SQL]
    * Check Convergence (Calculates the relative margins delta)
    bysort sex:  egen c_sex  = total(wt_raked)
    bysort area: egen c_area = total(wt_raked)
    
    gen d_sex  = abs(c_sex - cond(sex==1, `male_pop', `female_pop')) / cond(sex==1, `male_pop', `female_pop')
    gen d_area = abs(c_area - cond(area==1, `urban_pop', `rural_pop')) / cond(area==1, `urban_pop', `rural_pop')
    
    summarize d_sex, meanonly;  local max_d_sex = r(max)
    summarize d_area, meanonly; local max_d_area = r(max)
    local max_d = max(`max_d_sex', `max_d_area')
    
    drop c_sex c_area d_sex d_area
    if `max_d' < 0.0001 {
        local converged = 1
    }
}
\end{lstlisting}
\end{frame}

% --- SLIDE 25: POST-STRATIFICATION VS RAKING ---
\begin{frame}{Comparison: Post-Stratification vs. Raking}
  \begin{table}
    \centering
    \small
    \begin{tabular}{lll}
      \toprule
      \textbf{Feature} & \textbf{Post-Stratification} & \textbf{Raking} \\
      \midrule
      Number of variables & Single variable (or cross-classification) & Multiple variables simultaneously \\
      Benchmarks required & Joint distribution of variables & Marginal totals only \\
      Cell size requirement & Each cell needs sufficient sample ($n$) & Robust to smaller cell sizes \\
      Math complexity & Simple closed-form calculation & Iterative algorithm (converges) \\
      \bottomrule
    \end{tabular}
    \caption{Post-Stratification vs. Raking Methodology (Table 11)}
  \end{table}

  \begin{alertblock}{Key Thing to Hold Onto}
    Raking is the standard calibration method for large household surveys because it requires only marginal totals, making it highly practical to align multiple characteristics at once without cell size failures.
  \end{alertblock}
\end{frame}

% --- SECTION 6 ---
{
\setbeamercolor{background canvas}{bg=nsoblue}
\begin{frame}[plain, c]
  \centering
  \color{white}
  \Huge \textbf{Section 6: Implementing Unified Workflows} \\
  \vspace{0.4cm}
  \Large \textit{Sequential Weight Pipelines in Stata and Excel}
\end{frame}
}

% --- SLIDE 26: SEQUENTIAL WEIGHT PIPELINE ---
\begin{frame}{The Sequential Weight Pipeline}
  \begin{equation}
    w_i^{final} = \underbrace{\frac{1}{\pi_i}}_{\text{Design Weight}} \times \underbrace{\frac{n_c}{r_c}}_{\text{NR Adjustment}} \times \underbrace{\frac{N_g^{census}}{\hat{N}_g^{survey}}}_{\text{Calibration Factor}}
  \end{equation}

  \begin{columns}[T]
    \begin{column}{0.48\textwidth}
      \begin{block}{Sequential Workflow Goals}
        \begin{itemize}
          \item Each step targets a different source of bias.
          \item Final weights must sum to $N$.
          \item Complete estimation requires correcting standard errors.
        \end{itemize}
      \end{block}
    \end{column}
    \begin{column}{0.48\textwidth}
      \begin{block}{Stata Weight Rules (Table 12)}
        \begin{itemize}
          \item \textbf{Use \texttt{pweight} (probability weight):} Corrects standard errors using Taylor Linearization.
          \item \textbf{Do NOT use:} \texttt{fweight} (frequency), \texttt{aweight} (analytic), or \texttt{iweight} (importance). They distort standard errors.
        \end{itemize}
      \end{block}
    \end{column}
  \end{columns}
\end{frame}

% --- SLIDE 27: COMPLETE STATA SCRIPT (PART 1) ---
\begin{frame}[fragile]{Complete Pipeline in Stata (Part 1: Steps 1 \& 2)}
  This script takes raw survey data, builds design weights, and applies non-response cells:
  \begin{lstlisting}[language=SQL]
* Load raw survey data (includes respondents and non-respondents)
use "raw_survey_data.dta", clear

* STEP 1: Design Weight Construction
gen pi_overall = prob_psu * prob_hh_given_psu
gen wt_design = 1 / pi_overall

* STEP 2: Non-Response (NR) Adjustment
bysort stratum: egen n_sampled = count(hh_id)
bysort stratum: egen n_responded = total(responded)
gen nr_factor = n_sampled / n_responded
gen wt_nr = wt_design * nr_factor

* Retain responding units ONLY for subsequent analysis
keep if responded == 1
\end{lstlisting}
\end{frame}

% --- SLIDE 28: COMPLETE STATA SCRIPT (PART 2) ---
\begin{frame}[fragile]{Complete Pipeline in Stata (Part 2: Steps 3 to 5)}
  This completes the pipeline by setting post-stratification, declaring complex design, and analyzing:
  \begin{lstlisting}[language=SQL]
* STEP 3: Post-Stratification Alignment
gen ps_group = (sex - 1)*2 + area  // Group codes 1 to 4
bysort ps_group: egen survey_total = total(wt_nr)

gen census_pop = .
replace census_pop = 1850000 if ps_group == 1
replace census_pop = 1920000 if ps_group == 2
replace census_pop = 4100000 if ps_group == 3
replace census_pop = 4300000 if ps_group == 4

gen ps_factor = census_pop / survey_total
gen wt_final = wt_nr * ps_factor

* STEP 4: Declare Complex Survey Design to Stata
svyset psu [pweight = wt_final], strata(stratum) vce(linearized)

* STEP 5: Weighted National Estimation
svy: mean income
svy: total poor
\end{lstlisting}
\end{frame}

% --- SLIDE 29: EXCEL SPREADSHEET STRUCTURE ---
\begin{frame}[fragile]{Complete Excel Workbook Architecture}
  For training and quick point estimates, structure the workbook as follows:
  \begin{itemize}
    \item \textbf{Sheets:} \texttt{HH\_Data}, \texttt{Strata\_NR}, \texttt{PS\_Factors}, \texttt{Verification}, \texttt{Analysis}.
  \end{itemize}

  \begin{columns}[T]
    \begin{column}{0.48\textwidth}
      \begin{block}{\texttt{HH\_Data} Formulas (Listing 14)}
        \begin{itemize}
          \item \textbf{Overall Prob (Col F):} \\ \texttt{=prob\_psu * prob\_hh}
          \item \textbf{Design Weight (Col G):} \\ \texttt{=1/F2}
          \item \textbf{NR Factor Lookup (Col I):} \\ \texttt{=VLOOKUP(stratum, Strata\_NR!\$A:\$C, 3, FALSE)}
          \item \textbf{NR Adjusted Weight (Col J):} \\ \texttt{=G2 * I2}
          \item \textbf{Final Weight (Col M):} \\ \texttt{=J2 * L2} \textit{(PS\_Factor Applied)}
        \end{itemize}
      \end{block}
    \end{column}
    \begin{column}{0.48\textwidth}
      \begin{block}{Estimation Formulas (Listing 15)}
        \textbf{Weighted Mean Income:}
        \begin{lstlisting}[backgroundcolor=\color{white}, frame=none, basicstyle=\ttfamily\tiny]
=SUMPRODUCT(wt_final_range, income_range) 
 / SUM(wt_final_range)
        \end{lstlisting}
        \textbf{Weighted Poverty Rate:}
        \begin{lstlisting}[backgroundcolor=\color{white}, frame=none, basicstyle=\ttfamily\tiny]
=SUMPRODUCT(wt_final_range, poor_range) 
 / SUM(wt_final_range)
        \end{lstlisting}
        \textbf{Subpopulation (Urban) Mean Income:}
        \begin{lstlisting}[backgroundcolor=\color{white}, frame=none, basicstyle=\ttfamily\tiny]
=SUMPRODUCT((area_range=1)*wt_final_range, 
            income_range) 
 / SUMPRODUCT((area_range=1)*wt_final_range)
        \end{lstlisting}
      \end{block}
    \end{column}
  \end{columns}
\end{frame}

% --- SECTION 7 ---
{
\setbeamercolor{background canvas}{bg=nsoblue}
\begin{frame}[plain, c]
  \centering
  \color{white}
  \Huge \textbf{Section 7: Diagnostics and Common Mistakes} \\
  \vspace{0.4cm}
  \Large \textit{Validation Rules, Trimming Methods, and Kish's Design Effects}
\end{frame}
}

% --- SLIDE 30: CORE DIAGNOSTICS & WEIGHT RATIO ---
\begin{frame}{Core Diagnostics \& The Weight Efficiency Ratio}
  Before publishing any weighted estimates, run these key tests:
  \begin{columns}[T]
    \begin{column}{0.48\textwidth}
      \begin{block}{Diagnostic 1: Sum of Weights Check}
        The sum of final weights must equal target population size ($N$):
        \begin{itemize}
          \item \textbf{Stata:} \texttt{summarize wt\_final} $\Rightarrow$ verify \texttt{r(sum)}.
          \item \textbf{Excel:} \texttt{=SUM(wt\_final\_range)}.
        \end{itemize}
      \end{block}
    \end{column}
    \begin{column}{0.48\textwidth}
      \begin{block}{Diagnostic 2: The Weight Ratio}
        Extreme weights destabilize estimates and inflate variance. Evaluate the ratio:
        \begin{equation}
          \text{Weight Ratio} = \frac{w_{max}}{\overline{w}}
        \end{equation}
      \end{block}
    \end{column}
  \end{columns}

  \begin{table}
    \centering
    \footnotesize
    \begin{tabular}{ll}
      \toprule
      \textbf{Weight Ratio} & \textbf{Evaluation \& Guidelines} \\
      \midrule
      $< 3$ & Good; weights are well-behaved. \\
      $3 - 5$ & Acceptable; monitor estimates carefully. \\
      $5 - 10$ & Concerning; investigate causes of extreme weights. \\
      $> 10$ & Problematic; consider applying weight trimming. \\
      \bottomrule
    \end{tabular}
    \caption{Weight Ratio Evaluation Guidelines (Table 14)}
  \end{table}
\end{frame}

% --- SLIDE 31: WEIGHT TRIMMING ---
\begin{frame}[fragile]{Diagnostic 3: Weight Trimming Methods}
  Weight trimming caps extreme weights at a threshold and redistributes the excess weight to preserve population sum. This manages the \textbf{bias-variance tradeoff} (small bias added; variance reduced).

  \begin{columns}[T]
    \begin{column}{0.48\textwidth}
      \begin{block}{Method A: Percentile Trimming (Listing 16)}
\begin{lstlisting}[language=SQL]
* Cap weights at the 99th percentile
_pctile wt_final, p(99)
local_trim_threshold = r(r1)
gen wt_trimmed = min(wt_final, `trim_threshold')

* Renormalize to maintain population N
summarize wt_final, meanonly
local pop_total = r(sum)
summarize wt_trimmed, meanonly
replace wt_trimmed = wt_trimmed * (`pop_total' / r(sum))
\end{lstlisting}
      \end{block}
    \end{column}
    \begin{column}{0.48\textwidth}
      \begin{block}{Method B: Fixed Multiple Trimming (Listing 17)}
\begin{lstlisting}[language=SQL]
* Cap weights at k * mean (e.g. k=5)
summarize wt_final, meanonly
local mean_w = r(mean)
local k = 5
gen wt_trimmed = min(wt_final, `k' * `mean_w')

* Renormalize wt_trimmed to match pop_total
* (Follow renormalization logic from Method A)
\end{lstlisting}
      \end{block}
    \end{column}
  \end{columns}
\end{frame}

% --- SLIDE 32: ESS AND DEFF ---
\begin{frame}[fragile]{Diagnostic 4: Effective Sample Size \& Design Effect}
  \begin{block}{Effective Sample Size ($n_{eff}$)}
    The sample size of an equal-probability design (SRS) that would yield the exact same precision as the complex weighted sample:
    \begin{equation}
      n_{eff} = \frac{(\sum w_{i})^{2}}{\sum w_{i}^{2}}
    \end{equation}
  \end{block}
  \begin{block}{Design Effect ($DEFF_{weight}$)}
    Kish's approximate design effect due to weight variation ($CV_w$ is the coefficient of variation of the weights):
    \begin{equation}
      DEFF_{weight} = \frac{n}{n_{eff}} \approx 1 + CV_{w}^{2}
    \end{equation}
    \textit{A DEFF of 2.0 indicates that weight variation has doubled the variance (equivalent to halving your sample size).}
  \end{block}

  \begin{block}{Calculating Kish's ESS in Stata (Listing 18)}
\begin{lstlisting}[language=SQL]
gen wt_sq = wt_final^2
summarize wt_final, meanonly; local sum_w = r(sum)
summarize wt_sq, meanonly;    local sum_w_sq = r(sum)
local n_eff = (`sum_w'^2) / `sum_w_sq'
display "Effective Sample Size (ESS): " `n_eff'
display "Kish Weighting DEFF: " _N / `n_eff'
\end{lstlisting}
  \end{block}
\end{frame}

% --- SLIDE 33: BENCHMARK TABLE & EIGHT MISTAKES ---
\begin{frame}{Verification Tables and the Eight Common Mistakes}
  Always construct a \textbf{Benchmark Verification Table} (Table 15) to verify calibrated margins match population targets ($0.0\%$ difference). Avoid these common structural errors:
  
  \begin{columns}[T]
    \begin{column}{0.48\textwidth}
      \begin{block}{Common Mistakes (1 - 4)}
        \begin{enumerate}
          \item \textbf{Incorrect Weight Types:} Using frequency (\texttt{fweight}) or analytic weights instead of \texttt{pweights}.
          \item \textbf{Outdated Design Declarations:} Forgetting to update \texttt{svyset} when weight variables change.
          \item \textbf{Retaining Non-Respondents:} Failing to drop units where \texttt{responded == 0}.
          \item \textbf{Regression Over-Control:} Over-controlling for post-stratified variables without adjusting.
        \end{enumerate}
      \end{block}
    \end{column}
    \begin{column}{0.48\textwidth}
      \begin{block}{Common Mistakes (5 - 8)}
        \begin{enumerate}
          \item \textbf{Excel Row Duplication:} Treating weights as frequency numbers by duplicating rows.
          \item \textbf{Ignoring Weight Distribution:} Failing to inspect maximum weights for outlier errors.
          \item \textbf{Skipping Post-PS \texttt{svyset}:} Applying PS adjustments but calculating SEs with old weights.
          \item \textbf{Scale Mismatches:} Using normalized instead of unnormalized weights for aggregate totals.
        \end{enumerate}
      \end{block}
    \end{column}
  \end{columns}
\end{frame}

% --- SLIDE 34: PRE-PUBLICATION CHECKLIST (PART 1) ---
\begin{frame}{Pre-Publication Checklist: Data Validation}
  Before publishing survey results, check off every single step in the pipeline:
  \begin{block}{Weights \& Benchmark Verifications}
    \begin{itemize}
      \item[\checkmark] \textbf{Sum of Weights:} Confirm that the sum of the final weights matches the known target population size ($N$).
      \item[\checkmark] \textbf{Weight Ratios:} Ensure the weight ratio ($w_{max}/\overline{w}$) is less than 5.0, or weight trimming has been applied and documented.
      \item[\checkmark] \textbf{Weight Variation:} Verify the coefficient of variation ($CV_w$) and Effective Sample Size (ESS) have been calculated and reported.
      \item[\checkmark] \textbf{Benchmark Margins:} Construct a benchmark verification table to confirm all calibrated margins match targets.
    \end{itemize}
  \end{block}
\end{frame}

% --- SLIDE 35: PRE-PUBLICATION CHECKLIST (PART 2) ---
\begin{frame}{Pre-Publication Checklist: Code Execution}
  Confirm the software execution commands meet NSO methodologies:
  \begin{block}{Design Declarations \& Standard Errors}
    \begin{itemize}
      \item[\checkmark] \textbf{Survey Declaration:} Verify the survey design is correctly declared in Stata using \texttt{svyset} with the final weight variable, strata, and PSU variables.
      \item[\checkmark] \textbf{Estimation Prefix:} All point and variance estimates are produced using the \texttt{svy:} prefix.
      \item[\checkmark] \textbf{Exclusion of Non-Respondents:} Non-responding units have been dropped or excluded from the active analysis dataset.
      \item[\checkmark] \textbf{Methodology Documentation:} Trimming thresholds, weight types (normalized vs. unnormalized), and cell collapse assumptions are documented in the metadata.
    \end{itemize}
  \end{block}
  \begin{alertblock}{Key Summary}
    Survey weights are only as reliable as the diagnostics you run on them. Run systematic checks, document modifications, and verify final margins against census benchmarks before releasing national statistics.
  \end{alertblock}
\end{frame}

\end{document}